\chapter{完整 D-模}

设 $X$ 是光滑代数簇。以 $T^*(X)$ 记其余切丛，$p_X:T^*(X)\to X$ 记相应的自然投影。对任意凝聚
$\mathcal{D}_X$-模 $\mathcal{V}$，其特征簇 $\operatorname{Ch}(\mathcal{V})$ 是 $T^*(X)$ 的闭子簇。在映射 $p_X$ 下，特征簇
$\operatorname{Ch}(\mathcal{V})$ 投射到 $\mathcal{V}$ 的支集 $\operatorname{supp}(\mathcal{V})$ 上。我们已证明：对支集
$\operatorname{supp}(\mathcal{V})$ 中任意一点 $x$，有
$\dim_{p_X^{-1}(x)}\operatorname{Ch}(\mathcal{V})\ge\dim_xX$。因此，称凝聚 $\mathcal{D}_X$-模是\emph{完整（holonomic）的}，若对任意 $x\in\operatorname{supp}(\mathcal{V})$ 有
$\dim_{p_X^{-1}(x)}\operatorname{Ch}(\mathcal{V})=\dim_xX$。这推广了前述定义。

以 $\operatorname{Hol}(\mathcal{D}_X)$ 记 $\mathcal{M}_{coh}(\mathcal{D}_X)$ 中由所有完整 $\mathcal{D}_X$-模构成的全子范畴。

\section{完整 D-模}

\begin{theorem}\label{thm:5.1.1}
设 $X$ 是光滑簇。设
$$0\longrightarrow\mathcal{V}_1\longrightarrow\mathcal{V}_2\longrightarrow\mathcal{V}_3\longrightarrow0$$
是凝聚 $\mathcal{D}_X$-模的短正合列。则：
\begin{enumerate}[label=(\roman*)]
  \item 若 $\mathcal{V}_2$ 是完整模，则 $\mathcal{V}_1$ 与 $\mathcal{V}_3$ 也是完整的；
  \item 若 $\mathcal{V}_1$ 与 $\mathcal{V}_3$ 是完整的，则 $\mathcal{V}_2$ 是完整的。
\end{enumerate}
\end{theorem}

\begin{proof}
若 $\mathcal{V}_1$ 或 $\mathcal{V}_3$ 为零，断言显然成立。若它们非零，由前述可见
$\operatorname{Ch}(\mathcal{V}_2)=\operatorname{Ch}(\mathcal{V}_1)\cup\operatorname{Ch}(\mathcal{V}_3)$。这立即蕴含
$$\dim_{p_X^{-1}(x)}\operatorname{Ch}(\mathcal{V}_2)
=\max\bigl(\dim_{p_X^{-1}(x)}\operatorname{Ch}(\mathcal{V}_1),
\dim_{p_X^{-1}(x)}\operatorname{Ch}(\mathcal{V}_3)\bigr).$$
此外，由前述知 $\operatorname{supp}(\mathcal{V}_2)=\operatorname{supp}(\mathcal{V}_1)\cup\operatorname{supp}(\mathcal{V}_3)$。因此断言立即成立。
\end{proof}

因此，$\mathcal{M}_{coh}(\mathcal{D}_X)$ 的全子范畴 $\operatorname{Hol}(\mathcal{D}_X)$ 是 Abel 的且是厚的。

\begin{lemma}\label{lem:5.1.2}
设 $X$ 是光滑簇，$Y$ 是其闭光滑子簇。设 $i:Y\to X$ 是自然包含。则对任意 $\mathcal{D}_Y$-模 $\mathcal{V}$，以下陈述等价：
\begin{enumerate}[label=(\roman*)]
  \item 模 $\mathcal{V}$ 是完整的；
  \item 模 $i_+(\mathcal{V})$ 是完整的。
\end{enumerate}
\end{lemma}

\begin{proof}
显然，由 Kashiwara 定理，$\mathcal{V}$ 为零当且仅当 $i_+(\mathcal{V})$ 为零。此外，由前述，对任意 $x\in\operatorname{supp}(\mathcal{V})$ 有
$$\operatorname{holdef}_x(\mathcal{V})=\operatorname{holdef}_x(i_+(\mathcal{V})).$$
故对任意 $x\in\operatorname{supp}(\mathcal{V})$，$\dim_{p_Y^{-1}(x)}\operatorname{Ch}(\mathcal{V})=\dim_xY$ 当且仅当
$\dim_{p_X^{-1}(x)}\operatorname{Ch}(i_+(\mathcal{V}))=\dim_xX$。
\end{proof}

下一个结果推广了前述结论。

\begin{theorem}\label{thm:5.1.3}
每个完整 $\mathcal{D}_X$-模都是有限长度的。
\end{theorem}

\begin{proof}
显然，完整模在开集上的限制仍是完整的。由于由 III.2.8，为有限长度是局部性质，只需对光滑仿射簇证明该断言。此时可设 $X$ 是仿射空间 $k^n$ 的光滑闭子簇。设 $i:X\to k^n$ 是自然包含。则由引理 \ref{lem:5.1.2}，对完整模 $\mathcal{V}$，模 $i_+(\mathcal{V})$ 是完整的。由 I.8.1 可知 $i_+(\mathcal{V})$ 是有限长度的。显然 $i_+(\mathcal{V})$ 的任意子模或商模的支集都在 $X$ 中。故 $i_+(\mathcal{V})$ 作为所有 D-模范畴中对象的长度，等于它作为所有支集在 $X$ 中的模构成的子范畴中对象的长度。由 Kashiwara 定理，这意味着 $\mathcal{V}$ 是有限长度的。
\end{proof}

\section{联络（Connections）}

设 $X$ 是特征零的代数闭域 $k$ 上的光滑代数簇，$\mathcal{D}_X$ 是 $X$ 上微分算子层。本节描述一些非常简单的完整 $\mathcal{D}_X$-模。

若 $\mathcal{D}_X$-模 $\mathcal{V}$ 作为 $\mathcal{O}_X$-模是凝聚的，则称之为\emph{联络}。

\begin{theorem}\label{thm:5.2.1}
设 $X$ 是连通光滑代数簇，$\mathcal{V}$ 是非零凝聚 $\mathcal{D}_X$-模。则以下条件等价：
\begin{enumerate}[label=(\roman*)]
  \item $\mathcal{V}$ 是联络；
  \item $\mathcal{V}$ 的特征簇 $\operatorname{Ch}(\mathcal{V})$ 包含在余切丛 $T^*(X)$ 的零截面中；
  \item $\mathcal{V}$ 的特征簇 $\operatorname{Ch}(\mathcal{V})$ 就是余切丛 $T^*(X)$ 的零截面；
  \item $\mathcal{V}$ 是有限秩局部自由 $\mathcal{O}_X$-模。
\end{enumerate}
\end{theorem}

\begin{proof}
显然 (iv)$\Rightarrow$(i) 与 (iii)$\Rightarrow$(ii)。

\textbf{(iv)$\Rightarrow$(iii)} 由于 $\mathcal{V}$ 是有限秩局部自由 $\mathcal{O}_X$-模，几何纤维
$T_x(\mathcal{V})=\mathcal{V}_x/\mathfrak{m}_x\mathcal{V}_x$ 的维数是局部常值。因为 $X$ 连通，这意味着它是非零常值，从而
$\operatorname{supp}(\mathcal{V})=X$。可在 $\mathcal{V}$ 上定义滤过 $F\mathcal{V}$：当 $p<0$ 时 $F_p\mathcal{V}=0$，当 $p\ge0$ 时 $F_p\mathcal{V}=\mathcal{V}$。这显然是 $\mathcal{V}$ 上的好滤过。分次模 $\operatorname{Gr}\mathcal{V}$ 除
$\operatorname{Gr}_0\mathcal{V}=\mathcal{V}$ 外所有齐次分量都为 0。因此 $\operatorname{Gr}\mathcal{V}$ 的零化子包含
$\bigoplus_{p=1}^{\infty}\operatorname{Gr}_p\mathcal{D}_X$，从而 $\mathcal{V}$ 的特征簇包含在 $T^*(X)$ 的零截面中。由于由 III.3.9 它投射到 $\operatorname{supp}(\mathcal{V})$ 上，得出 $\operatorname{Ch}(\mathcal{V})$ 等于 $T^*(X)$ 的零截面。

\textbf{(ii)$\Rightarrow$(i)} 该陈述是局部的，故由 II.2.10，可设 $X$ 是仿射的且有整体坐标系
$(f_1,f_2,\dots,f_n;D_1,D_2,\dots,D_n)$，使得 $(D^I;I\in\mathbb{Z}_+^n)$ 构成 $X$ 上微分算子自由 $R(X)$-模的一组基。设 $F\mathcal{V}$ 是 $\mathcal{V}$ 的好滤过，$J$ 是 $\Gamma(X,\operatorname{Gr}\mathcal{V})$ 在 $R(T^*(X))$ 中的零化子。则 $J$ 在 $T^*(X)$ 中的零点集包含在 $T^*(X)$ 的零截面中。由 Hilbert 零点定理，$J$ 的根包含由 $D_1,D_2,\dots,D_n$ 的主符号生成的理想。这意味着存在 $m\in\mathbb{Z}_+$ 使得 $D_1^m,D_2^m,\dots,D_n^m$ 的主符号零化 $\Gamma(X,\operatorname{Gr}\mathcal{V})$。此外，任意阶 $q\ge nm$ 的微分算子的 $q$ 次主符号零化 $\Gamma(X,\operatorname{Gr}\mathcal{V})$。由于 $\Gamma(X,\operatorname{Gr}\mathcal{V})$ 是有限生成 $\operatorname{Gr}\mathcal{D}_X$-模，这意味着对充分大的 $p\in\mathbb{Z}_+$ 有 $\operatorname{Gr}_p\mathcal{V}=0$。因此对充分大的 $p\in\mathbb{Z}_+$ 有 $F_p\mathcal{V}=\mathcal{V}$，即 $\mathcal{V}$ 是凝聚 $\mathcal{O}_X$-模。

\textbf{(i)$\Rightarrow$(iv)} 由于该陈述是局部的，可设 $X$ 是仿射且“小”的。设 $x\in X$。由于 $\mathcal{V}$ 是凝聚 $\mathcal{O}_X$-模，几何纤维
$T_x(\mathcal{V})=\mathcal{V}_x/\mathfrak{m}_x\mathcal{V}_x=\mathcal{O}_{X,x}/\mathfrak{m}_x\otimes_{\mathcal{O}_{X,x}}\mathcal{V}_x$
是 $k$ 上有限维向量空间。设 $s_1,s_2,\dots,s_q$ 是 $\mathcal{V}$ 的整体截面族，其像 $s_1(x),s_2(x),\dots,s_q(x)$ 在 $T_x(\mathcal{V})$ 中构成该向量空间的一组基。这些截面定义了自由 $\mathcal{O}_X$-模 $\mathcal{O}_X^q$ 到 $\mathcal{V}$ 的自然态射。记其像为 $\mathcal{U}$。则有凝聚 $\mathcal{O}_X$-模的正合列
$$\mathcal{O}_X^q\longrightarrow\mathcal{V}\longrightarrow\mathcal{V}/\mathcal{U}\longrightarrow0,$$
它导致正合列
$$T_x(\mathcal{O}_X^q)\longrightarrow T_x(\mathcal{V})\longrightarrow T_x(\mathcal{V}/\mathcal{U})\longrightarrow0.$$
由于由构造第一个箭头是满射，得出 $T_x(\mathcal{V}/\mathcal{U})=0$，由 Nakayama 引理（I.2.1）得 $(\mathcal{V}/\mathcal{U})_x=0$。因此 $x$ 不在凝聚 $\mathcal{O}_X$-模 $\mathcal{V}/\mathcal{U}$ 的支集中。缩小 $X$ 可设 $\mathcal{V}/\mathcal{U}=0$，即 $\mathcal{V}$ 由 $s_1,s_2,\dots,s_q$ 生成。此时对 $X$ 上任意整体向量场 $T$ 有
$$Ts_i=\sum_{j=1}^q a^i_js_j$$
其中 $a^i_j\in R(X)$。欲证 $s_1,s_2,\dots,s_q$ 是自由 $\mathcal{O}_{X,x}$-模 $\mathcal{V}_x$ 的一组基。设 $\sum_i g_is_i=0$，其中 $g_1,\dots,g_q\in\mathcal{O}_{X,x}$。断言这意味着对任意 $p\in\mathbb{Z}_+$ 有 $g_1,\dots,g_q\in\mathfrak{m}_x^p$。

由于 $s_1(x),\dots,s_q(x)$ 线性无关，$\sum_i g_i(x)s_i(x)=0$ 蕴含
$g_1(x)=\dots=g_q(x)=0$，从而 $g_1,\dots,g_q\in\mathfrak{m}_x$。因此该陈述对 $p=1$ 成立。设它对 $p-1$ 成立。对 $X$ 上任意整体向量场 $T$ 有
$$0=T\sum_{i=1}^qg_is_i=\sum_{i=1}^q(T(g_i)s_i+g_iTs_i)$$
$$=\sum_{i=1}^qT(g_i)s_i+\sum_{i,j=1}^q a_j^ig_is_j
=\sum_{i=1}^q\Bigl(T(g_i)+\sum_{j=1}^qa_i^jg_j\Bigr)s_i.$$
故由归纳假设，$T(g_i)+\sum_{j=1}^qa_i^jg_j\in\mathfrak{m}_x^{p-1}$ 且 $g_i\in\mathfrak{m}_x^{p-1}$（$1\le i\le q$）。这意味着 $T(g_i)\in\mathfrak{m}_x^{p-1}$（$1\le i\le q$）。特别地，$D_j(g_i)\in\mathfrak{m}_x^{p-1}$（$1\le i\le q$，$1\le j\le n$）。这导致对所有 $|I|<p$ 的 $I\in\mathbb{Z}_+^n$ 有 $D^Ig_i\in\mathfrak{m}_x$，由 II.2.15 得出 $g_i\in\mathfrak{m}_x^p$（$1\le i\le q$）。

因此由对 $p$ 的归纳，对所有 $p\in\mathbb{Z}_+$ 有 $g_i\in\mathfrak{m}_x^p$（$1\le i\le q$）。故 $g_1=\dots=g_q=0$。于是 $s_1,\dots,s_q$ 是自由 $\mathcal{O}_{X,x}$-模 $\mathcal{V}_x$ 的一组基。

因此可考虑自然短正合列
$$0\longrightarrow\mathcal{K}\longrightarrow\mathcal{O}_X^q\longrightarrow\mathcal{V}\longrightarrow0,$$
其中 $\mathcal{K}$ 是自然态射 $\mathcal{O}_X^q\to\mathcal{V}$ 的核。显然 $\mathcal{K}$ 是凝聚 $\mathcal{O}_X$-模，由上述结果 $\mathcal{K}_x=0$。于是 $x\notin\operatorname{supp}(\mathcal{K})$，必要时缩小 $X$ 可设 $\mathcal{K}=0$。这意味着 $s_1,\dots,s_q$ 是自由 $\mathcal{O}_X$-模 $\mathcal{V}$ 的一组基。
\end{proof}

这个结果有以下明显推论。

\begin{corollary}\label{cor:5.2.2}
联络是完整 $\mathcal{D}_X$-模。
\end{corollary}

上述证明表明，若 $X$ 足够“小”且 $\mathcal{V}$ 是 $X$ 上联络，可找到自由 $\mathcal{O}_X$-模 $\mathcal{V}$ 的一组基 $s_1,\dots,s_q$，把 $\mathcal{V}$ 与 $\mathcal{O}_X^q$ 等同。设 $\Psi:\mathcal{O}_X^q\to\mathcal{V}$ 是由
$\Psi(f_1,\dots,f_q)=\sum_i f_is_i$ 给出的相应 $\mathcal{O}_X$-模同构。则存在 $q\times q$ 矩阵 $A_i=(A^i_{kj})$（$1\le i\le n$），其元素在 $R(X)$ 中，使得
$$D_is_j=\sum_{k=1}^qA^i_{kj}s_k.$$
于是
$$(D_i\circ\Psi)(f_1,\dots,f_q)=D_i\Bigl(\sum_{j=1}^qf_js_j\Bigr)=\sum_{j=1}^q\bigl(D_i(f_j)s_j+f_jD_is_j\bigr)$$
$$=\sum_{j=1}^qD_i(f_j)s_j+\sum_{j,k=1}^qf_jA^i_{kj}s_k
=\sum_{j=1}^q\Bigl(D_i(f_j)+\sum_{k=1}^qf_kA^i_{jk}\Bigr)s_j.$$
故
$$\Psi^{-1}\circ D_i\circ\Psi=D_i+A_i$$
对任意 $1\le i\le n$ 成立。此外，由于在 $\mathcal{V}$ 上 $[D_i,D_j]=0$，有
$$0=[D_i+A_i,D_j+A_j]=[D_i,A_j]+[A_i,D_j]+[A_i,A_j]=D_i(A_j)-D_j(A_i)+[A_i,A_j],$$
即
$$D_j(A_i)-D_i(A_j)=[A_i,A_j]$$
对所有 $1\le i,j\le n$ 成立。因此，联络局部地对应于可积联络（integrable connections）的经典概念。
\begin{annote}
$\mathcal O_X$-模配平坦联络 $\nabla$ 等价于局部系统（局部常值层）；取 de Rham 复形得到 $\mathrm{dR}:D^b_{rh}(X)\cong D^b_c(X^{an},\mathbb C)$（Mebkhout–Kashiwara–Beilinson–Bernstein）。这把“完整模的分类”翻译成“构造可层/反常层的分类”，是第五章所有结论的收束点。[来源：nLab 条目 flat+connection（Last revised 2023-06-04）、Riemann-Hilbert+correspondence（2024-08-09）；原书 §2, p133–138]
\end{annote}



\section{完整性与直像}

本节证明完整性在直像下保持。

先考虑光滑仿射簇态射的情形。设 $X$、$Y$ 是两个光滑仿射簇，$\varphi:X\to Y$ 是态射。以 $D_X$、$D_Y$ 分别记 $X$、$Y$ 上微分算子环，以 $\mathcal{M}_R(D_X)$、$\mathcal{M}_R(D_Y)$ 记相应的右 D-模范畴。则可考虑直像函子
$\varphi_+:\mathcal{M}_R(D_X)\to\mathcal{M}_R(D_Y)$ 及其左导出函子 $L^p\varphi_+$。

\begin{theorem}\label{thm:5.3.1}
设 $\mathcal{V}$ 是完整右 $D_X$-模。则 $L^p\varphi_+(\mathcal{V})$ 是完整右 $D_Y$-模。
\end{theorem}

\begin{proof}
在仿射空间 $X=k^n$、$Y=k^m$ 的情形下，这个结果已在 I.13.5 中证明。现在把定理的证明归结到该情形。显然可把 $X$、$Y$ 作为闭代数集嵌入仿射空间 $k^n$、$k^m$。设
$i_X:X\to k^n$、$i_Y:Y\to k^m$。由 Kashiwara 定理，直像函子
$i_{X,+}:\mathcal{M}^R_{qc}(D_X)\to\mathcal{M}^R_{qc}(D_{k^n})$ 与
$i_{Y,+}:\mathcal{M}^R_{qc}(D_Y)\to\mathcal{M}^R_{qc}(D_{k^m})$ 是正合的。借由记号上的滥用，用相同字母表示 $\mathcal{M}_R(D_X)$（相应地 $\mathcal{M}_R(D_Y)$）与 $\mathcal{M}_R(D^{(n)})$（相应地 $\mathcal{M}_R(D^{(m)})$）之间的相应直像函子。于是对 $X$ 有正合函子的交换图
$$\begin{tikzcd}
\mathcal{M}_R(D_X) \arrow[r,"i_{X,+}"] \arrow[d,"\operatorname{loc}"' ] &
\mathcal{M}_R(D^{(n)}) \arrow[d,"\operatorname{loc}"]\\
\mathcal{M}^R_{qc}(D_X) \arrow[r,"i_{X,+}"] & \mathcal{M}^R_{qc}(D_{k^n})
\end{tikzcd}$$
其中竖直箭头是局部化函子，$Y$ 有类似的图。由于竖直箭头是范畴的等价，而下箭头由 Kashiwara 定理是完全忠实的，得出上水平箭头也是完全忠实的，并建立 $\mathcal{M}_R(D_X)$ 与 $\mathcal{M}_R(D^{(n)})$ 中由支集在 $X$ 中的模构成的全子范畴的等价。

设 $\Phi:k^n\to k^m$ 是诱导 $\varphi:X\to Y$ 的多项式映射。则有态射的交换图
$$\begin{tikzcd}
X \arrow[r,"i_X"] \arrow[d,"\varphi"'] & k^n \arrow[d,"\Phi"]\\
Y \arrow[r,"i_Y"'] & k^m
\end{tikzcd}$$
由 IV.2.7 知，对任意 $p\in-\mathbb{Z}_+$ 有
$$L^p\Phi_+\circ i_{X,+}=L^p(\Phi_+\circ i_{X,+})=L^p(i_{Y,+}\circ\varphi_+)=i_{Y,+}\circ L^p\varphi_+.$$
设 $\mathcal{V}$ 是 $X$ 上完整右 $D_X$-模。由于由引理 \ref{lem:5.1.2}，函子 $i_{X,+}$ 把完整模映为完整模，模 $i_{X,+}(\mathcal{V})$ 是完整右 $D^{(n)}$-模。故由 I.13.5，$L^p\Phi_+(i_{X,+}(\mathcal{V}))$ 对所有 $p\in-\mathbb{Z}_+$ 是完整 $D^{(n)}$-模。这意味着 $i_{Y,+}(L^p\varphi_+(\mathcal{V}))$ 是完整 $D^{(n)}$-模。再次使用引理 \ref{lem:5.1.2}，得出对所有 $p\in-\mathbb{Z}_+$，$L^p\varphi_+(\mathcal{V})$ 是完整 $D_Y$-模。
\end{proof}

对定理 \ref{thm:5.3.1} 局部化可见：对任意完整右 $D_X$-模 $\mathcal{V}$，对任意 $p\in\mathbb{Z}$，右 $D_Y$-模 $H^p(\varphi_+(\mathcal{D}(\mathcal{V})))$ 是完整的。

现考虑一般情形。若 $\mathcal{D}_X$-模的复形 $\mathcal{V}^\cdot$ 满足所有 $p\in\mathbb{Z}$ 时 $H^p(\mathcal{V}^\cdot)$ 是完整 $\mathcal{D}_X$-模，则称它是\emph{完整复形}。以 $D^b(\mathcal{M}(\mathcal{D}_X))$ 记 $\mathcal{D}_X$-模的有界导出范畴。由于范畴 $\operatorname{Hol}(\mathcal{D}_X)$ 是 $\mathcal{D}_X$-模范畴的厚 Abel 子范畴，$D^b(\mathcal{M}(\mathcal{D}_X))$ 中由所有完整复形构成的全子范畴
$D^b_{hol}(\mathcal{M}(\mathcal{D}_X))$ 是三角子范畴。下一个结果是 [2] 的一个特例。

\begin{proposition}\label{prop:5.3.2}
完整 $\mathcal{D}_X$-模构成三角范畴 $D^b_{hol}(\mathcal{M}(\mathcal{D}_X))$ 的一个生成类。
\end{proposition}

事实上，可在 $D^b_{hol}(\mathcal{M}(\mathcal{D}_X))$ 中找到一个更小的生成类。设 $U\subset X$ 是仿射开集，$i:U\to X$ 是自然包含。设 $\mathcal{V}$ 是 $U$ 上完整模。断言 $i_+(\mathcal{V})$ 是 $X$ 上完整模。事实上，由于完整性是局部性质，只需证明对任意仿射开集 $V\subset X$，限制 $i_\bullet(\mathcal{V})|_V$ 是完整的。设 $j:U\cap V\to V$ 是自然包含。则
$i_\bullet(\mathcal{V})|_V=j_\bullet(\mathcal{V}|_{U\cap V})$，且由于 $U\cap V$ 仿射，由讨论的第一部分，$j_\bullet(\mathcal{V}|_{U\cap V})$ 是完整模。这证明了 $i_\bullet(\mathcal{V})$ 是完整模。

\begin{proposition}\label{prop:5.3.3}
设 $\mathcal{U}=(U_1,U_2,\dots,U_n)$ 是 $X$ 的仿射开覆盖，$i_j:U_j\to X$ 是自然包含。则模族 $\mathcal{G}$（由 $i_{j,\bullet}(\mathcal{V})$ 组成，其中 $\mathcal{V}$ 是 $U_j$ 上任意完整模且 $1\le j\le n$）构成
$D^b_{hol}(\mathcal{M}(\mathcal{D}_X))$ 的一个生成族。
\end{proposition}

\begin{proof}
设 $\mathcal{V}$ 是 $X$ 上完整模，$C^\cdot(\mathcal{U},\mathcal{V})$ 是其 Čech 分解。则 $\mathcal{V}$ 拟同构于
$C^\cdot(\mathcal{U},\mathcal{V})$。由相关结果，$\mathcal{G}$ 是 $D^b_{hol}(\mathcal{M}(\mathcal{D}_X))$ 的一个生成类。
\end{proof}

这最终使我们能够证明 I.13.5 的如下推广。

\begin{theorem}\label{thm:5.3.4}
设 $X$、$Y$ 是光滑代数簇，$\varphi:X\to Y$ 是簇的态射。若 $\mathcal{V}^\cdot$ 是有界完整 $\mathcal{D}_X$-模复形，则
$\varphi_+(\mathcal{V}^\cdot)$ 是有界完整 $\mathcal{D}_Y$-模复形。
\end{theorem}

\begin{proof}
由于完整性是局部性质，可设 $Y$ 是光滑仿射簇。

设 $U$ 是 $X$ 中仿射开集，$i:U\to X$ 是自然浸入，$\mathcal{V}$ 是 $U$ 上完整模。则 $i_+(\mathcal{V})$ 是完整模。由于由 IV.5.7，$i_+(\mathcal{D}(\mathcal{V}))=\mathcal{D}(i_\bullet(\mathcal{V}))$，可见
$$\varphi_+(\mathcal{D}(i_\bullet(\mathcal{V})))=\varphi_+(i_+(\mathcal{D}(\mathcal{V})))
=(\varphi\circ i)_+(\mathcal{D}(\mathcal{V})).$$
由于 $\varphi\circ i:U\to Y$ 是仿射簇之间的态射，由定理 \ref{thm:5.3.1} 的局部化，对 $p\in\mathbb{Z}$，
$H^p(\varphi_+(\mathcal{D}(i_\bullet(\mathcal{V}))))$ 是 $Y$ 上完整模。因此，由命题 \ref{prop:5.3.3}，结论成立。
\end{proof}

\section{不可约完整模的分类}

现给出不可约完整 $\mathcal{D}_X$-模的分类。它基于以下结果。

\begin{lemma}\label{lem:5.4.1}
设 $U$ 是 $X$ 的开子集，$i:U\to X$ 是自然浸入，$\mathcal{V}$ 是不可约完整 $\mathcal{D}_U$-模。则
\begin{enumerate}[label=(\roman*)]
  \item $i_\bullet(\mathcal{V})$ 包含唯一的不可约 $\mathcal{D}_X$-子模 $\mathcal{W}$；
  \item $\mathcal{W}|_U=\mathcal{V}$。
\end{enumerate}
\end{lemma}

\begin{proof}
由定理 \ref{thm:5.3.4}，$i_\bullet(\mathcal{V})$ 是完整 $\mathcal{D}_X$-模。因此由定理 \ref{thm:5.1.3}，它是有限长度的。设 $\mathcal{W}$ 是 $i_\bullet(\mathcal{V})$ 的不可约 $\mathcal{D}_X$-子模。由于 $i_\bullet$ 是限制到 $U$ 的函子的右伴随，$\mathcal{W}$ 的限制 $\mathcal{W}|_U$ 非零，因此等于 $\mathcal{V}$。这意味着 $i_\bullet(\mathcal{V})$ 的任意两个不可约 $\mathcal{D}_X$-子模的交非零，即 $\mathcal{W}$ 是 $i_\bullet(\mathcal{V})$ 的唯一不可约 $\mathcal{D}_X$-子模。
\end{proof}

这个结果蕴含不可约完整模的如下延拓结果。

\begin{corollary}\label{cor:5.4.2}
设 $U$ 是 $X$ 的开子集，$\mathcal{V}$ 是不可约完整 $\mathcal{D}_U$-模。则存在不可约完整 $\mathcal{D}_X$-模 $\mathcal{W}$，使得 $\mathcal{W}|_U$ 同构于 $\mathcal{V}$。此外，$\mathcal{W}$ 在同构意义下唯一，且 $\operatorname{supp}(\mathcal{W})$ 是 $\operatorname{supp}(\mathcal{V})$ 在 $X$ 中的闭包。
\end{corollary}

\begin{proof}
存在性部分立即由引理 \ref{lem:5.4.1} 得出。设 $\mathcal{W}'$ 是另一个不可约完整 $\mathcal{D}_X$-模，使得 $\mathcal{W}'|_U$ 同构于 $\mathcal{V}$。由于 $i_\bullet$ 是限制到 $U$ 的右伴随，存在自然态射
$\alpha:\mathcal{W}'\to i_\bullet(\mathcal{V})$，使得 $\alpha|_U$ 是 $\mathcal{W}'|_U$ 到 $\mathcal{V}$ 的同构。由于 $\mathcal{W}'$ 不可约，$\alpha$ 的核为零，其像是 $i_\bullet(\mathcal{V})$ 的不可约 $\mathcal{D}_X$-子模。由引理 \ref{lem:5.4.1}，它必等于 $\mathcal{W}$。

显然 $\operatorname{supp}(\mathcal{W})\cap U=\operatorname{supp}(\mathcal{V})$。因此 $\mathcal{V}$ 的支集在 $X$ 中的闭包包含在 $\mathcal{W}$ 的支集中。另一方面，$\mathcal{W}$ 的支集包含在 $\operatorname{supp}(i_\bullet(\mathcal{V}))$ 中，后者等于 $\operatorname{supp}(\mathcal{V})$ 在 $X$ 中的闭包。
\end{proof}

设 $V$ 是 $X$ 中连通光滑子簇，$j:V\to X$ 是自然浸入。设 $\tau$ 是 $V$ 上不可约联络。则直像
$$I(V,\tau)=j_+(\tau)$$
是完整 $\mathcal{D}_X$-模，称为附属于 $(V,\tau)$ 的\emph{标准 D-模}。由推论 \ref{cor:5.4.2}，其支集等于 $\bar{V}$。由于 $V$ 在 $X$ 中局部闭，可找到开集 $U\subset X$ 使得 $V$ 是 $U$ 的闭子簇。以 $j_U$ 记 $V$ 到 $U$ 的浸入，$i$ 记 $U$ 到 $X$ 的开浸入。则由前述有
$$I(V,\tau)=j_+(\tau)=i_+((j_U)_+(\tau)).$$
由 Kashiwara 定理，$(j_U)_+(\tau)$ 的直像是支集等于 $V$ 的不可约完整 $\mathcal{D}_U$-模。由引理 \ref{lem:5.4.1}，它延拓为不可约完整 $\mathcal{D}_X$-模，即 $I(V,\tau)$ 的唯一不可约子模。记为 $L(V,\tau)$，称为附属于数据 $(V,\tau)$ 的\emph{不可约模}。由推论 \ref{cor:5.4.2}，$L(V,\tau)$ 的支集等于 $\bar{V}$。因此我们证明了以下结果的第一部分。

\begin{proposition}\label{prop:5.4.3}
设 $I(V,\tau)$ 是附属于 $(V,\tau)$ 的标准 $\mathcal{D}_X$-模。则它包含唯一不可约子模 $L(V,\tau)$。$L(V,\tau)$ 的支集等于 $\bar{V}$，且 $I(V,\tau)/L(V,\tau)$ 的支集包含在 $\bar{V}-V$ 中。
\end{proposition}

\begin{proof}
只需证明 $\mathcal{Q}=I(V,\tau)/L(V,\tau)$ 的支集在 $\bar{V}-V$ 中。考虑短正合列
$$0\longrightarrow L(V,\tau)\longrightarrow I(V,\tau)\longrightarrow\mathcal{Q}\longrightarrow0.$$
令 $U=X-(\bar{V}-V)$。则 $U$ 是 $X$ 的开子集且 $V$ 在 $U$ 中闭。因此，由上述讨论与引理 \ref{lem:5.4.1}，从这个短正合列限制到 $U$ 得出 $\mathcal{Q}|_U=0$。故 $\operatorname{supp}(\mathcal{Q})\subset\bar{V}-V$。
\end{proof}

现在可以分类 $X$ 上不可约完整模。

\begin{theorem}\label{thm:5.4.4}
\begin{enumerate}[label=(\roman*)]
  \item 设 $\mathcal{V}$ 是不可约完整 $\mathcal{D}_X$-模。则存在 $\mathcal{V}$ 的支集中的不可约开光滑仿射子簇 $V$ 及 $V$ 上不可约联络 $\tau$，使得 $\mathcal{V}$ 同构于 $L(V,\tau)$。
  \item 设 $V$、$V'$ 是 $X$ 的两个不可约光滑仿射子簇，$\tau$、$\tau'$ 分别是 $V$、$V'$ 上不可约联络。则 $L(V,\tau)$ 同构于 $L(V',\tau')$ 当且仅当
    \begin{enumerate}[label=(\alph*)]
      \item $\bar{V}=\bar{V}'$；
      \item 存在 $V\cap V'$ 的非空开仿射子簇 $V''$ 使得 $\tau|_{V''}\cong\tau'|_{V''}$。
    \end{enumerate}
\end{enumerate}
\end{theorem}

\begin{proof}
\begin{enumerate}[label=(\roman*)]
  \item 由前述，$\mathcal{V}$ 的支集是 $X$ 的不可约闭子集。故存在 $X$ 的开仿射子集 $U$，使得 $V=\operatorname{supp}(\mathcal{V})\cap U$ 是 $U$ 的闭光滑子簇。显然 $\mathcal{V}|_U$ 是不可约完整 $\mathcal{D}_U$-模。若以 $j_V$ 记 $V$ 到 $U$ 的自然浸入，由 Kashiwara 定理，存在不可约完整 $\mathcal{D}_V$-模 $\mathcal{W}$，使得 $(j_V)_+(\mathcal{W})\cong\mathcal{V}|_U$。此外，$\operatorname{supp}(\mathcal{W})=V$。因此，由前述，存在 $V$ 的开稠密仿射子簇 $V'$ 使得 $\mathcal{W}|_{V'}$ 是联络。故必要时缩小 $U$，可进一步设 $\mathcal{W}$ 是联络。令 $\tau=\mathcal{W}$，则 $\mathcal{V}|_U$ 与 $L(V,\tau)|_U$ 同构。由推论 \ref{cor:5.4.2}，这意味着 $\mathcal{V}$ 与 $L(V,\tau)$ 同构。
  \item 若 $L(V,\tau)$ 同构于 $L(V',\tau')$，则它们的支集相等，(a) 成立。因此可设 $L(V,\tau)$ 与 $L(V',\tau')$ 有公共支集 $S$。于是 $V$ 与 $V'$ 在 $S$ 中开稠密，故 $V\cap V'$ 是 $V$ 与 $V'$ 的非空开仿射子簇。设 $V''$ 是 $V\cap V'$ 的非空开仿射子簇，则 $V''$ 不可约。设 $U$ 是 $X$ 的开子集使得 $V''=S\cap U$。则由推论 \ref{cor:5.4.2}，$L(V,\tau)|_U$ 与 $L(V',\tau')|_U$ 是支集等于 $V''$ 的同构的不可约完整 $\mathcal{D}_U$-模，当且仅当 $L(V,\tau)$ 与 $L(V',\tau')$ 同构。此外，若以 $j$ 记 $V''$ 到 $U$ 的浸入，则 $L(V,\tau)|_U=j_+(\tau|_{V''})$，$L(V',\tau')|_U=j_+(\tau'|_{V''})$。由于 $V''$ 是 $U$ 的光滑闭子簇，由 Kashiwara 定理，$j_+(\tau|_{V''})$ 同构于 $j_+(\tau'|_{V''})$ 当且仅当 $\tau|_{V''}$ 同构于 $\tau'|_{V''}$。
\end{enumerate}
\end{proof}
\begin{annote}
$\dim X=1$ 时完整模就是 $X$ 上有限秩平坦联络的推广；正则奇点情形下局部单值化给出 Deligne 的“正则奇点方程必有解”，即 $D_{rh}(X)$ 的每个对象都来自局部系统。这是全书中唯一能完全算出分类的情形，也是理解“完整亏格”的几何意义的入口。[来源：nLab 条目 Riemann-Hilbert+correspondence（2024-08-09）；原书 §2、§4, p133–142]
\end{annote}



\section{完整模的局部上同调}

设 $Y$ 是 $X$ 的闭光滑子簇，$U=X-Y$。以 $i:Y\to X$、$j:U\to X$ 记典范包含映射。

设 $\mathcal{V}^\cdot$ 是有界完整 $\mathcal{D}_X$-模复形。则 $\mathcal{V}^\cdot$ 到开集 $U$ 的限制 $\mathcal{V}^\cdot|_U$ 显然是 $\mathcal{D}_U$-模的有界完整复形。因此，由前述，$j_+(\mathcal{V}^\cdot|_U)=Rj_\bullet(\mathcal{V}^\cdot|_U)$ 是 $\mathcal{D}_X$-模的有界完整复形。由可区别三角
$$\begin{tikzcd}
& Rj_\bullet(\mathcal{V}^\cdot|_U) \arrow[dl,"+1"'] & \\
R\Gamma_Y(\mathcal{V}^\cdot) \arrow[rr] && \mathcal{V}^\cdot \arrow[lu]
\end{tikzcd}$$
可见 $R\Gamma_Y(\mathcal{V}^\cdot)$ 也是 $\mathcal{D}_X$-模的有界完整复形。

因此我们证明了以下结果。

\begin{lemma}\label{lem:5.5.1}
设 $\mathcal{V}^\cdot$ 是有界完整 $\mathcal{D}_X$-模复形。则 $R\Gamma_Y(\mathcal{V}^\cdot)$ 是有界完整 $\mathcal{D}_X$-模复形。
\end{lemma}

\section{完整性与逆像}

由前述结果已知 $R\Gamma_Y(\mathcal{V}^\cdot)=i_+(Ri^!(\mathcal{V}^\cdot))$。

应用 $Ri^!$ 并使用前述，可见 $Ri^!(\mathcal{V}^\cdot)$ 也是 $\mathcal{D}_Y$-模的有界完整复形。这进而蕴含 $Li_+(\mathcal{V}^\cdot)$ 是 $\mathcal{D}_Y$-模的有界完整复形。

因此我们证明了以下结果。

\begin{lemma}\label{lem:5.6.1}
设 $\mathcal{V}^\cdot$ 是有界完整 $\mathcal{D}_X$-模复形。则 $Li_+(\mathcal{V}^\cdot)$ 是有界完整 $\mathcal{D}_Y$-模复形。
\end{lemma}

另一方面，若 $p:X\times Y\to Y$ 是投影映射，则 $p_+(\mathcal{V}^\cdot)=\mathcal{O}_X\boxtimes\mathcal{V}^\cdot$。因此，$p_+:\mathcal{M}_{qc}(\mathcal{D}_Y)\to\mathcal{M}_{qc}(\mathcal{D}_{X\times Y})$ 是正合函子。此外有如下结果。

\begin{lemma}\label{lem:5.6.2}
设 $\mathcal{V}^\cdot$ 是有界完整 $\mathcal{D}_Y$-模复形。则 $p_+(\mathcal{V}^\cdot)$ 是有界完整 $\mathcal{D}_{X\times Y}$-模复形。
\end{lemma}

把这些结果与图构造相结合并使用前述，得到以下推广前述结论的结果。

\begin{theorem}\label{thm:5.6.3}
设 $f:X\to Y$ 是光滑代数簇的态射。则逆像函子 $Lf^+$ 把 $\mathcal{D}_Y$-模的有界完整复形映为 $\mathcal{D}_X$-模的有界完整复形。
\end{theorem}
